\documentclass[tikz,border=10pt]{standalone}
\usepackage{fontspec}
\setmainfont{DejaVu Serif}
\setsansfont{DejaVu Sans}
\setmonofont{DejaVu Sans Mono}
\usetikzlibrary{arrows.meta,positioning,shapes.geometric}
\tikzset{
  process/.style={rectangle,draw,rounded corners,align=center,
    text width=10cm,minimum height=10mm,font=\small},
  decision/.style={diamond,draw,aspect=3,align=center,
    text width=4cm,inner sep=1pt,font=\small},
  terminal/.style={rectangle,draw,rounded corners=5mm,
    align=center,text width=7cm,minimum height=8mm,font=\small},
  flow/.style={-{Stealth},thick}
}
\begin{document}
\begin{tikzpicture}[node distance=6mm and 18mm]
\node[terminal] (start) {Начало};
\node[process,below=of start] (n1) {\texttt{experiments/train.py}\\ Прочитать конфигурации; создать среду};
\draw[flow] (start) -- (n1);
\node[process,below=of n1] (n2) {\texttt{agents/dqn.py}\\ Создать основную и целевую сети из agents/network.py};
\draw[flow] (n1) -- (n2);
\node[process,below=of n2] (n3) {\texttt{agents/replay\_buffer.py}\\ Создать буфер опыта};
\draw[flow] (n2) -- (n3);
\node[process,below=of n3] (n4) {\texttt{environment/env.py}\\ Выполнить reset()};
\draw[flow] (n3) -- (n4);
\node[process,below=of n4] (n5) {\texttt{agents/dqn.py}\\ Выбрать случайное действие или лучшее по Q-оценкам};
\draw[flow] (n4) -- (n5);
\node[process,below=of n5] (n6) {\texttt{environment/env.py}\\ Выполнить step(action)};
\draw[flow] (n5) -- (n6);
\node[process,below=of n6] (n7) {\texttt{agents/replay\_buffer.py}\\ Сохранить переход};
\draw[flow] (n6) -- (n7);
\node[process,below=of n7] (n8) {\texttt{agents/dqn.py}\\ Если опыт накоплен: выбрать пакет и обучить сеть; обновить целевую сеть по расписанию};
\draw[flow] (n7) -- (n8);
\node[process,below=of n8] (n9) {\texttt{experiments/metrics.py}\\ Обновить показатели};
\draw[flow] (n8) -- (n9);
\node[process,below=of n9] (n10) {\texttt{experiments/train.py}\\ По расписанию сохранить метрики и состояние агента в runs/};
\draw[flow] (n9) -- (n10);

\node[decision,below=of n10] (budget) {Бюджет исчерпан?};
\draw[flow] (n10) -- (budget);
\node[terminal,below=of budget] (end) {\texttt{experiments/train.py}\\ Сохранить итог; закрыть среду};
\draw[flow] (budget) -- node[right] {Да} (end);
\node[decision,right=of budget] (done) {Эпизод завершён?};
\draw[flow] (budget) -- node[above] {Нет} (done);
\draw[flow] (done.east) -- ++(1,0) |- node[pos=0.25,right] {Да} (n4.east);
\draw[flow] (done.north) |- node[pos=0.25,left] {Нет} (n5.east);

\end{tikzpicture}

\end{document}
